\documentclass[12pt]{article}
\usepackage[utf8]{inputenc}
\usepackage[T1]{fontenc}
\usepackage{hyperref}
\title{Western Yiddish Language Resources for K-12 Curriculum Development}
\author{Compiled by Claude Sonnet 5.5 (claude-sonnet-5-5) for Andr\'as Kornai}
\date{}
\begin{document}
\maketitle
\section{Introduction}
Western Yiddish (ISO 639-3 yih) is dormant. No community members use it as a first language, and it is not known to be taught in schools. It was the main spoken language of Ashkenazi Jews in what is now Germany, France, Switzerland, northern Italy and the Netherlands between the 14th and 19th centuries. Swiss and Alsatian varieties are described as moribund. Little Western-specific teaching material turned up in the search, and this report does not pad that out.
\section{Primary Educational Resources}
No school primers or textbooks specific to Western Yiddish were found. The Omniglot page lists general (Standard) Yiddish textbooks, such as Weinreich's \emph{College Yiddish}, which are not Western-specific: \url{https://www.omniglot.com/books/language/yiddish.htm}
\section{Linguistic and Reference Materials}
The Jewish Languages Project profile covers the history, writing system and status of Western Yiddish: \url{https://www.jewishlanguages.org/western-yiddish}

The Endangered Languages Project page for Western Yiddish links to a user-compiled list of grammar, vocabulary and pronunciation sources: \url{https://endangeredlanguages.com/lang/2990/guide/6641}

The Princeton guide to the Verterbukh Yiddish dictionary describes a dictionary that includes regional and dialectal variants alongside standard forms: \url{https://libguides.princeton.edu/az/verterbukh-yiddish-dictionary}

UCF Libraries has a guide to the Yiddish (Judeo-German) call-number range, which is useful for locating books: \url{https://guides.ucf.edu/cnpj/pj5061}
\section{Religious and Community Resources}
No Western Yiddish religious or community resources were found in the search.
\section{Digital and Audio Resources}
No Western Yiddish audio or digital resources were found in the search.
\section{Curriculum Development Context}
There is no living first-language community and no known school programme. Any curriculum would be a revival or heritage-study effort. It would rely on archival and academic sources, and on adapting general Yiddish materials with care. Further searching under dialect names (Alsatian, Swiss, Judeo-German) is needed.
\end{document}